% TOP_PROBLEM: {"id": 2721, "title": "Kirby Problem 1.62", "category_id": 7, "category": {"id": 7, "name": "topology", "display_name": "Topology", "description": "Properties preserved under continuous deformations.", "slug": "topology", "order_index": 7, "created_at": "2026-07-31T15:26:25.671Z"}, "status": "open"}
% FIRST_POSTED: null
% ENTRY_KIND: statement_only
% Imported from: batch11.tex; UnsolvedMath ID 2721
\documentclass[11pt]{article}
\usepackage[margin=25mm]{geometry}
\usepackage{fontspec}
\setmainfont{Latin Modern Roman}
\setsansfont{Latin Modern Sans}
\usepackage{amsmath,amssymb,mathtools}
\usepackage{unicode-math}
\setmathfont{Latin Modern Math}
\usepackage{xurl,hyperref}
\setcounter{secnumdepth}{0}
\setlength{\parindent}{0pt}
\setlength{\parskip}{5pt}
\setlength{\emergencystretch}{3em}
\newcommand{\sref}[2]{\hyperlink{source:#1:#2}{[S#2]}}
\newcommand{\cataloguescope}[1]{\paragraph{Recorded target.}#1}
\begin{document}
% UnsolvedMath ID 2721; source catalogue code KP-1.62
\setcounter{section}{2720}
\section{Kirby Problem 1.62}\label{2721}
{\small\sffamily UnsolvedMath ID 2721 · Topology}
\subsection{Definitions and mathematical statement}

\paragraph{Shared notation.}$\mathbb N=\{1,2,\ldots\}$, and $\mathbb Z,\mathbb Q,\mathbb R,\mathbb C$ denote the integers, rationals, reals and complex numbers. Any other conventions are specified locally.


Let $L$ be an $n$-component link in $S^3$, and let $F_n$ be the free group on $n$ generators. A boundary-link structure is a surjection $\phi:\pi_1(S^3\setminus L)\twoheadrightarrow F_n$ taking meridians to a free basis. It is good if its kernel is perfect, meaning $\ker\phi=[\ker\phi,\ker\phi]$. A link is topologically slice if its components bound disjoint locally flat disks $D_1,\ldots,D_n$ in $B^4$. It is freely topologically slice if, in addition, the disk-exterior fundamental group is $F_n$, freely generated by the disk meridians. The untwisted Whitehead double replaces each component by the Whitehead pattern with zero longitudinal twisting; each clasp can be chosen positive or negative. The Borromean rings are the standard three-component link with pairwise unlinked components but nontrivial total linking.
\paragraph{Statement recorded in the supplied source.}
Is every good boundary link topologically slice? Is every good boundary link freely topologically slice? In particular, is the untwisted Whitehead double of the Borromean rings topologically slice for every choice of the three clasp signs? \sref{2721}{2}

\subsection{Short English statement}
Must good boundary links bound locally flat disks, even with free disk-complement group? Test this in particular on all clasp choices for the Whitehead double of the Borromean rings.
\subsection{Sources}
\begin{description}
\item[\hypertarget{source:2721:1}{S1}] ULAM.AI and the UnsolvedMath Contributors, UnsolvedMath: A Curated Collection of Open Mathematics Problems, record 2721 (KP-1.62). CC BY 4.0. \url{https://huggingface.co/datasets/ulamai/UnsolvedMath}.
\item[\hypertarget{source:2721:2}{S2}] R. İnanç Baykur, Robion C. Kirby and Daniel Ruberman, eds., \emph{K3: A New Problem List in Low-Dimensional Topology}, Mathematical Surveys and Monographs 295, American Mathematical Society (2026), Problem 1.62 and its remarks; Chapters 1 and 2 for the stated conventions. \url{https://math.berkeley.edu/sites/default/files/surv-295-ruberman-watermarked-author-pdf.pdf}.
\end{description}
\label{end:2721}
\end{document}
